%*****************************************
\chapter{Percevoir l'environnement sonore}\label{ch:psycho_ea}
%*****************************************

Avant d'aller plus avant dans la présentation de nos recherches, ils nous est indispensable ici de dresser un état des lieux des connaissances liées à la perception des sons.

Nous proposons de présenter ces connaissances en quatre parties. La première récapitule les processus de traitement connus et mises en œuvre à partir du moment où le signal atteint le récepteur, autrement dit le tympan. 

Dans un second temps, nous présentons les résultats émanant de différentes études, regroupées a posteriori sous l'appellation Analyse de Scènes Acoustiques (ASA), et qui s’intéressent à la manière dont le cerveau assimile/ségrègue  les différentes informations contenues dans l'environnement sonore afin d'en dégager des objets cohérents, \ie des sources sonores. 

La troisième partie présente quant à elle les résultats d'études s'attachant a comprendre le fonctionnement  des mécanismes haut niveaux opérant lors de l’évaluation perceptive des environnements sonores, en empruntant une méthodologie issue de la psychologie cognitive. Il s'agit en particulier, d'étudier la contribution, sur l'évaluation globale, des objets composant ces environnements.

Enfin, la dernière partie résume, de manière non exhaustives, ce que le domaine des neurosciences nous apprend sur le traitement de l'information sonore par le cerveau. 

\section{Des oreilles à la pensée: Traitements cognitifs de l'information auditive}

\subsection{La cognition}

\subsection{Le traitement de l'information auditive}

Le son, comme nous l'entendons dans ce document, est une vibration émise par une source d'excitation, et transmise à l'air. Cette vibration se propage ensuite jusqu'à atteindre un récepteur, le tympan, qui va capter le différentielle de pression résultant de cette vibration. 

Comme décrit dans \cite{mcadams1994penser}, une fois le tympan atteint, le son passe par plusieurs processus de traitement, dont les principaux sont représentés Figure~\ref{fig:traitementSonMcAdamsBigand}.

\begin{figure}[bth]
        \myfloatalign
        \includegraphics[width=.5\linewidth]{gfx/traitementSonMcAdamsBigand}
        \caption[Principaux processus de traitement de l'information auditive et leurs interactions]{Principaux processus de traitement de l'information auditive et leurs interactions, à partir de \cite{mcadams1994penser}}\label{fig:traitementSonMcAdamsBigand}
\end{figure}

\begin{itemize}
\item \textit{Transduction}: 
\item \textit{Processus de groupement auditif}: (cf. \ref{sec:ASA})
\item \textit{Extraction/Calcul des propriétés/attributs}: 
\item \textit{Structure de connaissances abstraites}: 
\item \textit{Traitement des structure d'événements}: 
\item \textit{Représentation mentale de l'environnement sonore}: (cf. \ref{sec:psychoCog})
\end{itemize}

\subsection{Processus Bottum-up et processus Top-down}


\section{Analyse de scènes acoustiques}
\label{sec:ASA}

L'Analyse de Scène Acoustique (ASA) est champ de recherche qui s'est développé dans les années 1990, notamment grâce aux travaux de Albert S. Bregman \cite{bregman1994auditory}. L'objectif de l'ASA est de comprendre comment l'humain fait sens de l'agglomérat de sons distincts qui l'entourent. En particulier, il s'agit d'étudier la faculté du cerveau à isoler de ce magma sonore les informations relatives à une même source, et à les regrouper afin de les traiter comme un tout cohérent.

L'ASA a permis de montrer que la perception des sons suit, en grande partie, les mêmes règles que celles observés pour la vision, et notamment les principes décrits par la \textit{Gestalt Theory}.

\section{Évaluation et représentation mentale de l'environnement sonore}
\label{sec:psychoCog}

\subsection{psychologie cognitive}

\subsection{L'étude des paysages sonores}

\subsection{L'approche catégorielle}

Les objectifs de l'approche catégorielle sont doubles : il s'agit 1) d'appréhender les principes psychologiques qui sous-tendent la formation des représentations mentales évoquées à la section~\ref{fig:traitementSonMcAdamsBigand}, ainsi que 2) de comprendre l'influence de ces représentations sur le traitement de l'information sonore. 

\subsubsection{Théorie de la catégorisation}

La plupart des études adoptant cette approche s'appuie sur la théorie de la catégorisation comme initialement formalisée par E. Rosh et B. B. Lloyd \cite{roach1978cognition}. Il est à noter que l'étude de la catégorisation est antérieure aux travaux de Rosh, et que d'autres théories ont depuis été proposées. Mais les travaux de Rosh ayant été les premiers à présenter une formalisation stricte des principes régissant la catégorisation, et étant par ailleurs à l'origine des théories postérieures, nous nous appuierons sur ces derniers afin de détailler ces processus perceptifs. 


Un des principes essentiel de tout être vivant est de segmenter son environnement, \ie~de se bâtir un système de classification permettant de regrouper des objets n'étant pas identiques. On appelle catégorisation, l'action consistant à regrouper des objets du monde physique considérés comme équivalents, et catégorie, la classe contenant le groupe d'objets ainsi rassemblé. Il est important de noter qu'à chaque catégorie est associé un label, censé décrire du mieux possible l'ensemble des objets inclus dans la catégorie. On parlera alors de catégories sémantique. 

L'introduction de cette dimension sémantique a des conséquences importantes sur l'universalité de la catégorisation. Cette dernière étant codée par la langue, elle ne dépend plus seulement d'un contexte physique, mais également d'un contexte culturel. La catégorisation peut être vue comme un action intermédiaire entre, d'une part l'organisation d'une connaissance individuelle résultant d'une expérience sensorielle personnelle, et la constitution d'une représentation collective pouvant être partagée par le biais d'un langage commun \cite{dubois2006cognitive}. 

Selon Rosh la catégorisation doit obéir à deux principes \citep[p. 29]{roach1978cognition}:

\begin{enumerate}
\item \textit{L'économie cognitive}: la catégorisation doit fournir un maximum d'information pour un minimum d'effort. La structure du système catégorielle, \ie de la représentation mentale, doit s’élaborer en tenant compte de se principe d'économie, externe à la réalité physique des objets. Ainsi la catégorisation n'est pas absolue, mais elle répond à une demande, et est donc susceptible de varier suivant le contexte de cette demande.
\item \textit{La structure du monde perçu}: l'ensemble des objets physiques ne vie pas dans un espace aux dimensions bien identifiées, indépendantes et finies, et dont les valeurs seraient équiprobables. Comme énoncé par D. Dubois \citep[p. 33]{dubois1991semantique} ``\,il existerait dans le monde des discontinuités, ainsi que des corrélations ou co-occurences de propriétés (telles que, pour un oiseau, avoir à la fois des plumes, un bec, deux pattes ...). Le monde des objets naturels ne constitue ainsi qu'un ensemble des possibles [$\ldots$] \,''. L'organisation catégorielle de nos représentations mentales se structure en fonction de ces discontinuités et co-occurrences de propriétés. 
\end{enumerate}



Rosh propose une vision prototypique de la catégorisation, \ie~les catégories se structurent en interne autour de prototypes qui 

\subsubsection{Catégories de paysages sonores}

\subsection{L'approche dimensionnelle}

\section{Neuroscience}
\label{sec:Neuroscience}
TODO
%*****************************************
%*****************************************
%*****************************************
%*****************************************
%*****************************************
